\chapter{Ablauf}

Generell gilt:
\begin{itemize}
    \item Themen von Praxis- und Abschlussarbeiten sind mit einem  Betreuer der FHDW abzusprechen und im FHDW-Portal unter \url{https://my.fhdw-hannover.de/} anzumelden. Der erste Ansprechpartner für die Betreuung ist Ihr Mentor. 
    \item Arbeiten müssen vor dem Ablauf der Abgabefrist digital im Portal \url{https://my.fhdw-hannover.de/} unter \textit{Abgaben} hochgeladen werden. \textbf{Zur Abgabe gehören folgende Dokumente}: Diese sind als separate Dateien hochzuladen:
    \begin{enumerate}
        \item Das Hauptdokument der Arbeit als PDF, inklusive der Anhangs-Kapitel (Appendizes), sofern vorhanden.
        \item Die \textit{Ehrenwörtliche Erklärung}, eigenhändig  unterschrieben und eingescannt als PDF.
        \item Bei Bedarf weitere Dokumente wie Fragebögen, Rohdaten, Source\-code oder ähnliches als ZIP-Datei.
    \end{enumerate}
    \item In allen Fällen von Terminüberschreitungen wird die Arbeit mit \textit{nicht ausreichend} bewertet!
    \item Eine Praxisarbeit kann bei Nichtbestehen jeweils einmal wiederholt werden. Darüber hinaus gelten im Einzelnen die unten aufgeführten Regeln.
\end{itemize}

\section{Praxisprojekt}
\label{sec:praxisprojekt}

\begin{itemize}
    \item Ein obligatorisches Praxisprojekt ist im 2., 3. oder 4. Praxisquartal zu erstellen. Empfohlen wird die Anfertigung im 3. Praxisquartal.
    \item Die reine Bearbeitungszeit der Praxisprojekte beträgt \textbf{6 Wochen}, die sich beliebig innerhalb der 12 Wochen des Praxisquartals verteilen können. Das Kolloquium sollte incl.\ Vorbereitung maximal eine Woche Aufwand erzeugen. Für ein Praxisprojekt mit Kolloquium werden 11 CP bzw.\ 12 CP im Master vergeben.
    \item Das Thema kann in Ihrem Partnerunternehmen vergeben werden und ist rechtzeitig vor Bearbeitungsbeginn mit einem von Ihnen gewählten Betreuer abzusprechen. 
    \item Den Bearbeitungsbeginn und alle weiteren Termine entnehmen Sie dem jeweiligen Zeitplan. Eine Verlängerung des Abgabetermins ist nicht möglich. Im Falle von Terminüberschreitungen wird die Arbeit mit \textit{nicht ausreichend} bewertet.
    \item Mit dem FHDW-Betreuer werden im Allgemeinen zwei Termine während der Anfertigungszeit durchgeführt: Der erste Termin dient dazu, den geplanten Inhalt der Arbeit anhand eines Inhaltsverzeichnisses zu diskutieren und so zu sichern, dass die Arbeit fachgerecht aufgebaut ist. Vor dem zweiten Termin soll der Betreuer die Gelegenheit für eine Leseprobe erhalten. Dazu sollen zwischen 1/3 und 2/3 der Arbeit fertiggestellt sein. Im Termin selbst können die bisherigen Inhalte dann diskutiert werden. 
    \item Das Kolloquium findet in der Regel im folgenden Theoriequartal statt. Es soll die Praxisarbeit inhaltlich in einem 20-minütigen Vortrag mit 10 Minuten anschließender Diskussion vor Publikum präsentieren. 
\end{itemize}

\section{Bachelor-Thesis}

\begin{itemize}
    \item Die Bachelor-Thesis (Bachelor-Abschlussarbeit) wird in der Regel im 6. Praxisquartal angefertigt. Es werden incl.\ Kollquium 16 CP vergeben.
    
    \item Im Kolloquium verteidigen Sie die Inhalte Ihrer Thesis in einem Abschlussvortrag oder einem Prüfungsgespräch mit den beiden Gutachtern (s. auch Abschitt~\ref{sec:bewertung-kolloquium}). Der Erstprüfer legt die Modalität des Kolloquiums fest. Klären Sie die Modalität frühzeitig mit Ihrem Erstprüfer. 
    
    \item Die Bearbeitungszeit für die Abschlussarbeit beträgt \textbf{10 Wochen}. Die Termine entnehmen Sie bitte dem Zeitplan des jeweiligen Quartals. 
    
    \item \textbf{Tipp}: Zur Absprache der Themenstellung verfassen Sie vor dem Beginn der Bachelor-Thesis eine \textbf{Kurz-Exposé} des geplanten Inhalts inklusive eines groben Projektplans. Besprechen Sie dieses Kurz-Exposé mit Ihrem Betreuer und ggf. Ihrem Ansprechpartner im Partnerunternehmen. Sie können folgende \LaTeX{}-Vorlage dafür nutzen:
    
    \begin{center}
    \url{https://www.overleaf.com/read/hmpmfjprfzqf}.
    \end{center}

\end{itemize}


\section{Master-Thesis}

\begin{itemize}
    \item Die Master-Thesis (Master-Abschlussarbeit) wird im Anschluss an das 3. Mastersemester angefertigt. Es werden 20 CP vergeben. 
    \item Die Bearbeitungszeit für die Abschlussarbeit beträgt \textbf{16 Wochen}.
    \item Bzgl. Terminen, Verlängerungen, abzugebende Dokumente, Art des Kolloquiums gelten die selben Regeln wie bei der Bachelor-Thesis.
    \item \textbf{Tipp}: Insbesondere bei einer Master-Thesis raten wir Ihnen zur schriftlichen Absprache der Themenstellung mit Ihrem Erstprüfer und ggf. dem Praxisbetreuer mittels des \textbf{Kurz-Exposés} (siehe oben).
\end{itemize}

\section{Transferprojekt}

\begin{itemize}
    \item Das Transferprojekt ist Teil des Aufbaustudiums als Vorbereitung für einen Masterstudiengang. Es werden 10 CP vergeben.
    \item Weitere Details zum Ablauf besprechen Sie mit Ihrem Prüfer.
\end{itemize}


\section{Regelungen zu Krankheit und Verlängerung der Bearbeitungszeit aus triftigen Gründen}

Die genauen Regelungen definiert der allgemeine Teil der Prüfungsordnung (siehe InfoPool auf my-fhdw.de \url{https://my.fhdw-hannover.de/infopool.php?p_id=41&ipac=details&fileId=7551}).

\begin{itemize}
    \item \textbf{Im Krankheitsfall}:
    
    \textit{\enquote{Der Bearbeitungszeitraum von Praxis- und Abschlussarbeiten werden um Krankheitsphasen verlängert, wenn diese unverzüglich attestiert werden.}} (s. PO allgemeiner Teil) 
    
    Das bedeutet, dass Sie im Krankheitsfall unverzüglich ein Attest einreichen müssen, damit die Bearbeitungszeit verlängert werden kann. Reichen Sie das Attest erst dann nach, wenn Sie merken, dass Ihnen zum Ende der Bearbeitungszeit die Zeit fehlt, ist eine Verlängerung der Bearbeitungszeit nicht mehr möglich.
    
    \item \textbf{Andere triftige Gründe} 
    (gilt nur für \textit{Abschlussarbeiten (Bachelor- oder Master-Thesis)}): 
    
    \textit{\enquote{Der Prüfungsausschuss kann in antragsmäßig begründeten Einzelfällen die Bearbeitungszeit um drei Wochen verlängern.}} (s. PO allgemeiner Teil).

    Können Sie Ihre Abschlussarbeit aus begründeten und nicht selbstverschuldeten Gründen nicht fristgemäß fertigstellen, kann ein Antrag auf die Verlängerung der Bearbeitungszeit um bis zu drei Wochen gestellt werden.
    Begründete und nicht selbstverschuldete Gründe sind Gründe, die nicht allein auf eine schlechte Zeitplanung oder mangelnde Arbeitsdisziplin des Studierenden zurückzuführen sind.

    Ein entsprechender Antrag muss \textit{3 Wochen vor Ende der regulären Bearbeitungszeit} in Absprache mit dem Erstgutachter gestellt werden. Den genauen spätesten Zeitpunkt hierfür finden Sie im Intranet der FHDW Hannover in den Informationen zu Bachelor- und Master-Thesis.
\end{itemize}

